\chapter{Document Extraction and Ingestion Pipeline}
\label{ch:extraction}

This chapter specifies the complete end-to-end pipeline for converting regulatory
filings and company disclosures into structured, source-linked observations that
AccDSL can reason about. It addresses the gap in v1.2 where extraction was
mentioned but not designed.

\section{Pipeline overview and quality contract}

\subsection{The extraction challenge}

Financial disclosures are published in diverse formats:
\begin{itemize}
\item \textbf{XBRL/iXBRL:} Structured but complex (contexts, dimensions, calculation
  linkbases, extension taxonomies)
\item \textbf{PDF:} Unstructured (multi-column layouts, tables spanning pages,
  embedded images, scanned documents)
\item \textbf{HTML:} Semi-structured (CSS layouts, JavaScript-rendered tables, responsive designs)
\item \textbf{Word/Excel:} Native office formats (requires conversion)
\end{itemize}

Extraction must preserve:
\begin{itemize}
\item \textbf{Source fidelity:} Original bytes, page/element locators, literal text
\item \textbf{Structure:} Tables, hierarchies, cross-references, footnotes
\item \textbf{Semantics:} Concepts, units, periods, dimensions, sign conventions
\item \textbf{Provenance:} Extraction method, confidence, review decisions
\end{itemize}

\subsection{Quality contract}

The extraction pipeline must achieve:

\begin{longtable}{L{35mm}L{50mm}L{51mm}}
\toprule Quality dimension & Target & Measurement\\\midrule\endhead
\textbf{Document coverage} & 100\% of acquired pages/sections inventoried &
Count pages with disposition vs. total pages\\
\textbf{Source linking} & $\geq$99\% of structured observations have valid locator &
Spot checks + automated validation\\
\textbf{Table structure} & $\geq$95\% of financial tables correctly parsed
(row/column headers, cell alignment) & GriTS metric on annotated corpus\\
\textbf{Numeric accuracy} & $\geq$99.9\% of extracted numbers match source (no
transcription errors) & Digit-by-digit comparison\\
\textbf{Context correctness} & $\geq$95\% of observations have correct entity,
period, unit, currency & Human annotation on heldout set\\
\textbf{Concept mapping} & $\geq$85\% of line items correctly mapped to standard
concepts; $\geq$95\% recall for material accounts & Precision/recall vs. expert annotations\\
\textbf{Cross-reference resolution} & $\geq$90\% of note references resolved to
correct targets & Link-traversal validation\\
\textbf{Fabrication rate} & $\leq$0.1\% of extracted facts not in source &
False-positive audit\\
\bottomrule
\end{longtable}

Extraction below these thresholds for a given filing must be flagged for human review.

\section{Stage 1: Document acquisition and inventory}

\subsection{Filing manifest construction}

Input: Company identifier (CIK, LEI, ticker), filing date range, filing types
(10-K, 10-Q, 20-F, etc.).

\textbf{Process:}
\begin{enumerate}
\item Query EDGAR API or equivalent for matching filings
\item Download complete filing packages (HTML/XML, exhibits, inline XBRL)
\item Compute SHA-256 hash of each file
\item Extract metadata (accession number, filing date, period end, document type)
\item Store original bytes in content-addressed storage
\item Create filing manifest with document inventory
\end{enumerate}

\textbf{Manifest schema:}

\begin{lstlisting}
FilingManifest {
  filing_id: String,
  entity: EntityIdentifier,
  filing_date: Date,
  period_end: Date,
  filing_type: Enum[10K, 10Q, 8K, 20F, ...],
  documents: [
    {
      hash: SHA256,
      filename: String,
      media_type: String, // application/pdf, text/html, application/xml
      size_bytes: Integer,
      storage_uri: URI
    }
  ],
  exhibits: [ExhibitReference],
  incorporated_references: [IncorporatedReference],
  acquisition_timestamp: DateTime,
  acquisition_method: String
}
\end{lstlisting}

\subsection{Content structure discovery}

For each document, detect:
\begin{itemize}
\item \textbf{PDF:} Page count, TOC (if present), section headings via font size/style heuristics
\item \textbf{HTML:} DOM structure, `<h1>`--`<h6>` headings, `<table>` elements,
  XBRL inline facts (`<ix:nonFraction>`, `<ix:nonNumeric>`)
\item \textbf{XBRL:} Instance document, taxonomy schema, linkbases, calculation relationships
\end{itemize}

Output: Document structure map with section/note inventory and expected vs. actual content.

\begin{lstlisting}
DocumentStructure {
  document_hash: SHA256,
  content_units: [
    {
      unit_id: String,
      type: Enum[Cover, TOC, Section, Table, Paragraph, Exhibit],
      locator: PageNumber | ElementPath,
      parent_unit: String?,
      heading_text: String?,
      expected_in_manifest: Boolean,
      detected: Boolean
    }
  ],
  missing_units: [ExpectedUnit],
  unexpected_units: [DetectedUnit]
}
\end{lstlisting}

\section{Stage 2: Layout analysis and text extraction}

\subsection{PDF parsing with Docling}

Docling (MIT license, inspect model licenses separately) provides:
\begin{itemize}
\item Page layout analysis (identify text blocks, tables, figures, headers/footers)
\item Reading order determination
\item Table structure recognition (row/column spans, headers)
\item Text extraction with bounding boxes
\end{itemize}

\textbf{Integration contract:}

\begin{lstlisting}
def parse_pdf(pdf_path: Path, config: DoclingConfig) -> DoclingResult:
    """
    Returns:
      - pages: List[PageLayout] with text blocks and tables
      - extraction_quality: Dict[str, float] with confidence scores
      - errors: List[ExtractionError] for failed pages/tables
    """
    ...
\end{lstlisting}

\textbf{Configuration:}
\begin{itemize}
\item Table detection threshold (default 0.7)
\item OCR fallback for scanned pages (enable if text extraction empty)
\item Language hint (English default; affects OCR)
\item Preserve original page images for review
\end{itemize}

\textbf{Quality checks:}
\begin{enumerate}
\item \textbf{Text extraction completeness:} Compare character count with expected
  from file size; flag if $<50\%$ extracted (possible encryption or scanning)
\item \textbf{Table detection:} Validate detected table count vs. expected from
  financial-statement structure (balance sheet, income statement, cash flow $\geq 3$)
\item \textbf{Reading order:} Check that page $n$ content precedes page $n+1$
  content (detect multi-column reordering errors)
\end{enumerate}

\subsection{HTML parsing and inline XBRL extraction}

HTML filings (especially inline XBRL) embed structured facts in rendered text.

\textbf{Process:}
\begin{enumerate}
\item Parse HTML with BeautifulSoup or lxml
\item Extract inline XBRL facts: `<ix:nonFraction>`, `<ix:nonNumeric>`, `<ix:hidden>`
\item For each fact, extract:
  \begin{itemize}
  \item Context reference (entity, period, dimensions)
  \item Unit reference (currency, shares)
  \item Concept QName (namespace + local name)
  \item Literal value and format
  \item Decimal/precision attributes
  \item Parent element path (for source linking)
  \end{itemize}
\item Extract tables as structured data (row/column headers from `<th>`, data from `<td>`)
\item Identify cross-references (`<a href="#note5">`, resolve to target section)
\end{enumerate}

\textbf{Quality checks:}
\begin{enumerate}
\item \textbf{XBRL validation:} Use Arelle to validate instance against taxonomy;
  report calculation inconsistencies
\item \textbf{Context completeness:} Every fact must resolve to a valid context
  with entity, period, dimensions
\item \textbf{Unit completeness:} Every numeric fact must have a unit (USD, shares, pure)
\item \textbf{Duplicate facts:} Flag duplicate (context, concept) pairs; may indicate
  restatements or errors
\end{enumerate}

\section{Stage 3: Table structure recognition}

Financial tables are the most information-dense part of filings. Correct structure
is critical for numerical extraction.

\subsection{Table components}

A financial table consists of:
\begin{itemize}
\item \textbf{Caption:} Table title (e.g., "Note 5 – Debt")
\item \textbf{Column headers:} Periods, dimensions, currencies (often multi-level)
\item \textbf{Row headers:} Line item names (often hierarchical with indentation)
\item \textbf{Data cells:} Numeric values, sometimes with footnote markers
\item \textbf{Footnotes:} Cell-specific qualifications
\end{itemize}

\subsection{Recognition pipeline}

\textbf{Step 1: Detect table boundaries}

Use Docling table detector or Table Transformer (baseline from PubTables-1M).
Output: Bounding boxes for each table.

\textbf{Step 2: Identify header rows/columns}

Heuristics:
\begin{itemize}
\item Top rows with bold or different font: likely column headers
\item Left column with text while right columns numeric: likely row headers
\item Cells spanning multiple columns: higher-level headers
\item Repeated patterns (e.g., "Current Year" / "Prior Year"): period headers
\end{itemize}

Use GTE (financial table benchmark) trained model if available.

\textbf{Step 3: Parse hierarchical structure}

Row headers often have indentation indicating parent-child relationships:

\begin{verbatim}
Debt
  Current portion                    10
  Long-term debt
    Senior notes                     50
    Subordinated debt                20
  Total debt                         80
\end{verbatim}

Extract:
\begin{itemize}
\item Indentation level (leading spaces or tab count)
\item Parent-child relationships
\item Totals and subtotals (often bold or indented differently)
\end{itemize}

\textbf{Step 4: Canonicalize to structured representation}

Output schema:

\begin{lstlisting}
TableStructure {
  table_id: String,
  caption: String,
  source_locator: PageNumber | ElementPath,
  column_headers: [
    {
      level: Integer, // 0 = top level
      span: (col_start, col_end),
      text: String,
      dimension: Enum[Period, Segment, Currency, Category]
    }
  ],
  row_headers: [
    {
      row_index: Integer,
      text: String,
      indentation_level: Integer,
      parent_row: Integer?,
      is_total: Boolean
    }
  ],
  cells: [
    {
      row: Integer,
      col: Integer,
      value: String, // Raw text before parsing
      numeric_value: Decimal?,
      unit: String?, // e.g., "millions", "$"
      footnote_refs: [String]
    }
  ]
}
\end{lstlisting}

\subsection{Quality validation}

\begin{enumerate}
\item \textbf{Structural consistency:} Every cell must be in a valid (row, col) position
\item \textbf{Arithmetic checks:} If table has labeled totals, verify sum of components
  equals total (within rounding)
\item \textbf{Column alignment:} Numeric cells should be right-aligned; text left-aligned
\item \textbf{Duplicate headers:} Flag tables with identical column headers (possible parsing error)
\item \textbf{Missing cells:} Tables should be dense; many empty cells suggests
  structure mis-recognition
\end{enumerate}

Compare against GriTS metric (cell-level structure accuracy) on annotated validation set.

\section{Stage 4: Semantic mapping and concept resolution}

Raw tables and text must be mapped to accounting concepts, periods, and entities.

\subsection{Concept mapping}

\textbf{Challenge:} Company labels vary; standard concepts are fixed.

Examples:
\begin{itemize}
\item "Accounts Receivable, Net" $\rightarrow$ \code{us-gaap:AccountsReceivableNetCurrent}
\item "Trade and other payables" $\rightarrow$ \code{ifrs-full:TradeAndOtherCurrentPayables}
\item "Property, Plant \& Equipment" $\rightarrow$ \code{us-gaap:PropertyPlantAndEquipmentNet}
\end{itemize}

\textbf{Approach:}

\begin{enumerate}
\item \textbf{Exact match:} Check if company label matches a standard-concept label
\item \textbf{Taxonomy extension:} If XBRL, company may have defined extension
  concepts with explicit base relationships
\item \textbf{Embedding similarity:} Use pre-trained financial-domain language
  model (e.g., FinBERT) to compute semantic similarity between company label and
  standard concepts; rank candidates
\item \textbf{Hierarchical constraints:} Use calculation linkbase or table structure
  to constrain mappings (if "Total Assets" is sum of components, each component must
  be an asset subconcept)
\item \textbf{Human review:} Present top-3 candidates with confidence scores; analyst
  selects or provides alternative
\end{enumerate}

\textbf{Mapping record:}

\begin{lstlisting}
ConceptMapping {
  source_label: String,
  source_context: Context, // Entity, period, dimensions
  target_concept: QName,
  relationship: Enum[Equivalent, Component, Broader, Reclassification, Transformation],
  confidence: Float in [0,1],
  mapping_method: Enum[ExactMatch, TaxonomyExtension, Embedding, HumanReview],
  reviewer: String?,
  review_date: DateTime?,
  evidence: [SourceSpan]
}
\end{lstlisting>

\subsection{Period and dimension resolution}

\textbf{Period normalization:}

Financial filings use diverse period representations:
\begin{itemize}
\item "Year Ended December 31, 2024"
\item "Three Months Ended June 30, 2024"
\item "Fiscal Year 2024" (company-specific fiscal year)
\item "52 weeks ended September 28, 2024"
\end{itemize}

Normalize to:
\begin{lstlisting}
Period {
  calendar: Calendar, // Gregorian, Fiscal52Week, etc.
  type: Enum[Instant, Duration],
  start_date: Date?,
  end_date: Date, // Instant = observation date; Duration = exclusive end
}
\end{lstlisting}

\textbf{Dimension extraction:}

XBRL dimensions (segment, product, geography, etc.) must be preserved:

\begin{lstlisting}
DimensionValue {
  axis: QName, // e.g., us-gaap:StatementGeographicalAxis
  member: QName, // e.g., country:US
  source_label: String // "United States"
}
\end{lstlisting}

\subsection{Unit and currency parsing}

Extract from:
\begin{itemize}
\item XBRL unit references (straightforward)
\item Table headers: "In thousands", "In millions, except per share"
\item Currency symbols: \$, €, £, ¥
\item Footnotes: "All amounts in Canadian dollars unless stated"
\end{itemize}

Output:
\begin{lstlisting}
Unit {
  base_currency: ISO4217, // USD, EUR, GBP, JPY
  scale: Integer, // 0 = units, 3 = thousands, 6 = millions
  per_unit: Unit?, // For ratios: $/share, $/barrel
}
\end{lstlisting}

\section{Stage 5: Cross-reference traversal and note linking}

Financial statements reference notes; notes reference exhibits and other filings.

\subsection{Link extraction}

Identify references:
\begin{itemize}
\item Explicit hyperlinks: `<a href="#note5">` in HTML
\item Text patterns: "See Note 5", "refer to Note 12"
\item Exhibit references: "filed as Exhibit 10.1 to Form 8-K"
\item Prior filings: "incorporated by reference from Annual Report on Form 10-K
  filed March 15, 2024"
\end{itemize}

\textbf{Resolution process:}

\begin{enumerate}
\item Parse link text to extract target (note number, exhibit number, filing accession)
\item Search document structure for matching target
\item If target in current document: resolve to locator
\item If target in exhibit: check exhibit inventory
\item If target in prior filing: check acquisition status
\item Record resolution status: Resolved | Duplicate | Unavailable | Ambiguous
\end{enumerate}

\textbf{Link record:}

\begin{lstlisting}
CrossReference {
  source_locator: Locator,
  link_text: String,
  target_type: Enum[Note, Exhibit, PriorFiling, ExternalURL],
  target_identifier: String,
  target_locator: Locator?,
  resolution_status: Enum[Resolved, Duplicate, Unavailable, Ambiguous],
  traversal_depth: Integer // How many hops from root
}
\end{lstlisting}

\subsection{Traversal algorithm}

\begin{lstlisting}
def traverse_references(manifest: FilingManifest, max_depth: int = 3):
    queue = [(root_statement, 0) for root_statement in manifest.primary_statements]
    visited = set()
    unresolved = []

    while queue:
        locator, depth = queue.pop(0)
        if locator in visited or depth > max_depth:
            continue
        visited.add(locator)

        links = extract_links(locator)
        for link in links:
            target = resolve_link(link, manifest)
            if target.status == Resolved:
                queue.append((target.locator, depth + 1))
            else:
                unresolved.append((link, target.status))

    return {
        'visited': visited,
        'unresolved': unresolved,
        'coverage': len(visited) / expected_content_units(manifest)
    }
\end{lstlisting}

\section{Stage 6: Observation creation and validation}

Raw extraction output is converted to immutable observation records.

\subsection{Observation schema}

From Chapter 2 v1.2, expanded:

\begin{lstlisting}
Observation {
  observation_id: UUID,
  // Source provenance
  document_hash: SHA256,
  filing_id: String,
  source_locator: Locator, // Page + bbox or element path
  literal_text: String, // Original rendering
  extraction_version: String,
  extraction_timestamp: DateTime,

  // Reporting context
  entity: EntityIdentifier,
  perimeter: ConsolidationScope,
  concept: QName,
  concept_mapping: ConceptMapping,
  taxonomy_version: String,
  dimensions: [DimensionValue],

  // Amount and precision
  original_string: String,
  parsed_value: Decimal?,
  sign_convention: Enum[Debit, Credit, Absolute],
  unit: Unit,
  scale: Integer,
  rounding_interval: Decimal?,
  missing_category: Enum[NotDisclosed, NotApplicable, Unreadable, Withheld]?,

  // Temporal context
  period: Period,
  economic_validity: (Date, Date)?, // When known
  publication_date: Date,
  availability_date: DateTime,
  ingestion_timestamp: DateTime,

  // Quality and review
  extraction_confidence: Float in [0,1],
  validation_errors: [ValidationError],
  reviewer_decision: Enum[Proposed, Accepted, AcceptedAlternative, Rejected]?,
  reviewer: String?,
  review_date: DateTime?,
  review_reason: String?,

  // Relationships
  supersedes: [ObservationID],
  correction_of: ObservationID?,
  conflicts_with: [ObservationID],
  components: [ObservationID], // For totals
  related_narrative: [NarrativeID]
}
\end{lstlisting}

\subsection{Validation rules}

Before committing an observation:

\begin{enumerate}
\item \textbf{Source link valid:} Locator resolves to actual content in document
\item \textbf{Context complete:} Entity, concept, period all present; dimensions
  valid for concept's taxonomy
\item \textbf{Value parseable:} If numeric, string parses to valid decimal; sign
  convention applied correctly
\item \textbf{Unit compatible:} Monetary concepts have currency; share counts have
  "shares" unit; ratios have per-unit
\item \textbf{Period type match:} Balance-sheet items have instant periods; income
  items have duration
\item \textbf{Temporal consistency:} Publication date $\geq$ period end date;
  ingestion $\geq$ publication
\item \textbf{No self-contradiction:} If observation supersedes another, contexts
  must overlap; correction must change value
\end{enumerate}

Failed validations produce warnings or errors based on severity. Errors block
observation creation; warnings flag for review.

\subsection{Duplicate resolution}

Filings often contain the same fact multiple times (statement + note, original + restated).

\textbf{Deduplication strategy:}

\begin{enumerate}
\item Compute canonical key: (entity, concept, period, dimensions, unit)
\item Group observations by canonical key
\item Within group:
  \begin{itemize}
  \item If values identical (within rounding): Keep all source links but create
    single observation with multiple locators
  \item If values differ: Check if one is labeled "As Restated" or "As Corrected"
    $\rightarrow$ create supersession relationship
  \item If values differ without explanation: Flag as conflict; retain both with
    conflict link; require human review
  \end{itemize}
\end{enumerate}

\section{Stage 7: Review interface and human-in-the-loop}

Extraction cannot be fully automated; human review is required for:
\begin{itemize}
\item Low-confidence extractions
\item Concept-mapping ambiguity
\item Conflicting values
\item Narrative assertions and qualitative disclosures
\item Novel transaction types
\end{itemize}

\subsection{Review queue prioritization}

Not all extractions need immediate review. Prioritize by:

\begin{longtable}{L{40mm}L{40mm}L{56mm}}
\toprule Priority & Criteria & Rationale\\\midrule\endhead
\textbf{Critical (P0)} & Material account (assets, liabilities, equity, revenue,
profit); confidence $<0.8$; conflicts detected & Affects financial position materially\\
\textbf{High (P1)} & Non-material accounts with conflicts; concept mapping
confidence $<0.7$ & May affect ratios or reconciliations\\
\textbf{Medium (P2)} & Note details, supplementary disclosures; confidence $<0.6$ &
Needed for complete model, not immediate\\
\textbf{Low (P3)} & Boilerplate policies, narrative with no quantitative impact &
Can be deferred or omitted\\
\bottomrule
\end{longtable}

\subsection{Review UI requirements}

The interface must present:

\textbf{Left pane: Source document}
\begin{itemize}
\item Original PDF page or HTML rendering
\item Highlighted bounding box or element
\item Surrounding context (full table, adjacent paragraphs)
\end{itemize}

\textbf{Right pane: Proposed interpretation}
\begin{itemize}
\item Extracted value, unit, period, dimensions
\item Proposed concept mapping with confidence and alternatives
\item Validation warnings/errors
\item Related observations (duplicates, conflicts, components)
\end{itemize}

\textbf{Actions:}
\begin{itemize}
\item Accept: Mark as Accepted, proceed to next
\item Edit: Modify value, concept, period, or unit; re-validate
\item Alternative: Create alternative interpretation (multiple readings coexist)
\item Reject: Mark as extraction error; record reason
\item Skip: Defer to later review
\item Link: Connect to related observations or narrative
\end{itemize}

\textbf{Workflow:}
\begin{itemize}
\item Reviewer assigned queue based on domain expertise
\item Progress tracked: N accepted / M total in queue
\item Consensus protocol: Two reviewers for P0/P1; one for P2/P3
\item Disagreements escalated to senior accountant
\end{itemize}

\section{Stage 8: Reconciliation and completeness checks}

After observations are created and reviewed, validate against expected financial-statement structure.

\subsection{Statement completeness}

Expected statements for annual filing:
\begin{itemize}
\item Balance Sheet (at least 2 periods: current, prior)
\item Income Statement (at least 2 periods)
\item Statement of Comprehensive Income (or combined with income)
\item Cash Flow Statement (at least 2 periods)
\item Statement of Changes in Equity (at least 2 periods)
\end{itemize}

Check:
\begin{enumerate}
\item All expected statements present with observations
\item Required line items present (assets, liabilities, equity, revenue, expenses,
  profit, cash flows)
\item Periods consistent (if income statement is 2024 and 2023, balance sheet
  should be Dec 31 2024 and Dec 31 2023)
\end{enumerate}

\subsection{Arithmetic reconciliation}

\textbf{Balance sheet identity:}
\[
\text{Total Assets} = \text{Total Liabilities} + \text{Total Equity}
\]

Check both current and prior periods. Tolerate rounding differences (typically $\pm$1
in display units).

\textbf{Income reconciliation:}
\[
\text{Net Income} = \text{Revenue} - \text{Expenses} - \text{Taxes}
\]

\textbf{Cash flow reconciliation:}
\[
\text{Cash}_{\text{end}} = \text{Cash}_{\text{start}} + \text{Operating CF} +
\text{Investing CF} + \text{Financing CF} + \text{FX Effect}
\]

\textbf{Equity rollforward:}
\[
\text{Equity}_{\text{end}} = \text{Equity}_{\text{start}} + \text{Profit} +
\text{OCI} - \text{Dividends} + \text{Issuance} - \text{Repurchase}
\]

For each identity, report:
\begin{itemize}
\item Matched: Both sides equal (or within rounding tolerance)
\item Unmatched: Difference exceeds tolerance; require investigation
\item Incomplete: Missing observations for one or more terms
\end{itemize}

\subsection{Cross-note validation}

Notes must reconcile with statement amounts:
\begin{itemize}
\item Debt note detail should sum to debt on balance sheet
\item Segment revenue should sum to total revenue
\item Tax expense components should sum to total tax expense
\end{itemize}

Create validation rules for common reconciliations; flag exceptions for review.

\section{Performance targets and benchmarks}

\subsection{Throughput}

Target processing speed:
\begin{itemize}
\item \textbf{XBRL filing:} $<$5 minutes for full pipeline (fast-path; structured)
\item \textbf{PDF annual report (100-200 pages):} $<$30 minutes for extraction +
  validation (human review separate)
\item \textbf{Human review:} ~5-10 observations per minute for experienced reviewer
  (varies by complexity)
\end{itemize}

For corpus of 100 companies $\times$ 3 filings each = 300 filings:
\begin{itemize}
\item Automated extraction: ~50 hours compute time (parallelizable)
\item Human review (assume 20\% of observations need review, 500 obs/filing):
  $300 \times 500 \times 0.2 \div 5 \text{ obs/min} = 6{,}000$ minutes $\approx$ 100 hours
\end{itemize}

\subsection{Evaluation on benchmark datasets}

\textbf{PubTables-1M:} Table structure recognition
\begin{itemize}
\item Metric: GriTS (cell-level structure accuracy)
\item Baseline: Table Transformer (~0.85 GriTS on PubTables test set)
\item Target: $\geq$0.90 on financial-table subset
\end{itemize}

\textbf{FinTabNet (GTE):} Financial table understanding
\begin{itemize}
\item Metric: Row/column header F1, cell extraction accuracy
\item Baseline: GTE model (~0.75 F1)
\item Target: $\geq$0.85 F1 on standard table types
\end{itemize}

\textbf{TAT-QA:} Tabular reasoning with questions
\begin{itemize}
\item Metric: Exact match (EM) and F1 on numerical answers
\item Baseline: SOTA models ~40-50\% EM
\item Use case: Validate extracted observations can answer known questions about filing
\end{itemize}

\textbf{AccDSL custom corpus:} Annotated real filings
\begin{itemize}
\item 3 companies $\times$ 3 filings = 9 filings fully annotated
\item Metrics: Observation precision/recall, concept-mapping accuracy, reconciliation success rate
\item Target: $\geq$95\% recall on material accounts, $\geq$85\% precision on concepts
\end{itemize}

\section{Error budget and continuous improvement}

\subsection{Error taxonomy}

Classify extraction errors:

\begin{longtable}{L{30mm}L{50mm}L{56mm}}
\toprule Error type & Description & Mitigation\\\midrule\endhead
\textbf{Transcription} & Wrong digit extracted (7 vs. 1, OCR error) &
Confidence scoring, manual verification of low-confidence\\
\textbf{Structure} & Table rows/columns misaligned & Improved layout model, heuristic validation\\
\textbf{Context} & Wrong period, entity, or dimension assigned & NER + co-reference
resolution, human review\\
\textbf{Concept} & Wrong concept mapping & Better semantic embeddings, calculation-linkbase
constraints\\
\textbf{Fabrication} & Fact not in source (hallucination) & Source-link validation,
audit random sample\\
\textbf{Omission} & Fact in source but not extracted & Recall measurement, heuristic
expected-item checks\\\bottomrule
\end{longtable}

\subsection{Monitoring and feedback loop}

\begin{enumerate}
\item \textbf{Per-filing quality report:} After extraction, generate report with:
  \begin{itemize}
  \item Document coverage percentage
  \item Observation count by confidence tier
  \item Validation errors and warnings
  \item Reconciliation status
  \item Unresolved cross-references
  \end{itemize}

\item \textbf{Human reviewer feedback:} Track:
  \begin{itemize}
  \item Accept rate (should be $>$80\% for P1/P2 observations)
  \item Common rejection reasons
  \item Time per observation (identify bottlenecks)
  \end{itemize}

\item \textbf{Downstream validation:} When observations feed into AccDSL historical
  reconstruction:
  \begin{itemize}
  \item Infeasibility $\rightarrow$ extraction error (wrong value or context)
  \item Missing data $\rightarrow$ omission error
  \item Reconciliation failure $\rightarrow$ concept-mapping or arithmetic error
  \end{itemize}

\item \textbf{Retraining triggers:} Retrain extraction models when:
  \begin{itemize}
  \item Accept rate drops below 70\% on new filing batch
  \item New company/industry patterns emerge (e.g., first insurance filing)
  \item Standards change (new IFRS/GAAP requirements)
  \end{itemize}
\end{enumerate}

\subsection{Model improvement protocol}

\begin{enumerate}
\item Collect failed cases (rejected by human or downstream validation)
\item Annotate correct extraction for failed cases
\item Add to training set or few-shot examples
\item Retrain/fine-tune extraction models
\item Validate on held-out set
\item Deploy new model version with A/B test (route 10\% of filings to new model,
  compare accept rates)
\item If new model improves, promote to production
\end{enumerate}

\section{Implementation roadmap}

\subsection{Milestone breakdown}

\begin{longtable}{L{20mm}L{45mm}L{71mm}}
\toprule Milestone & Deliverable & Dependencies\\\midrule\endhead
\textbf{M1} & Document acquisition from EDGAR; manifest construction; SHA-256 hashing &
HTTP client, SEC EDGAR API\\
\textbf{M2} & PDF parsing with Docling; text extraction; table detection &
Docling installation, model downloads\\
\textbf{M3} & HTML + inline XBRL parsing with Arelle; fact extraction &
Arelle installation, taxonomy downloads\\
\textbf{M4} & Table structure recognition; row/column header detection &
PubTables model, heuristic rules\\
\textbf{M5} & Concept mapping with FinBERT embeddings; similarity search &
FinBERT model, taxonomy concept embeddings\\
\textbf{M6} & Observation schema implementation; validation rules; database storage &
SQLite, schema migrations\\
\textbf{M7} & Review UI (web-based); source + interpretation display; workflow &
React/Vue frontend, REST API\\
\textbf{M8} & Reconciliation engine; arithmetic checks; completeness validation &
Observation query API\\
\textbf{M9} & End-to-end test on Apple 10-K fixture; quality report &
All prior milestones\\
\textbf{M10} & Benchmark evaluation on PubTables, FinTabNet, TAT-QA &
Test harness, annotated datasets\\\bottomrule
\end{longtable}

\subsection{Estimated effort}

\begin{itemize}
\item \textbf{Software engineering:} 2 FTE $\times$ 6 months = 12 person-months
  (pipeline infrastructure, integration, UI)
\item \textbf{ML engineering:} 1 FTE $\times$ 4 months = 4 person-months (model
  selection, fine-tuning, evaluation)
\item \textbf{Domain expertise (accounting):} 0.5 FTE $\times$ 6 months = 3 person-months
  (concept-mapping rules, validation logic, review protocol)
\item \textbf{QA + annotation:} 2 FTE $\times$ 3 months = 6 person-months (corpus
  annotation, testing, human review)
\end{itemize}

Total: ~25 person-months for production-ready extraction pipeline.

\section{Integration with AccDSL core}

Extraction pipeline outputs feed directly into AccDSL evidence registry (Chapter 2 v1.2).

\textbf{Interface contract:}

\begin{lstlisting}
def ingest_filing(
    filing_manifest: FilingManifest,
    extraction_config: ExtractionConfig
) -> EvidenceVersion:
    """
    Returns:
      - evidence_version_id: Unique ID for this extraction
      - observations: List[Observation] with source links
      - validation_report: Quality metrics and errors
      - review_queue: Observations requiring human review
    """
    ...
\end{lstlisting}

AccDSL compiler (Chapter 6 v1.2) queries evidence registry for observations matching
requested entity, period, concepts. Extraction pipeline ensures these observations
are source-linked, validated, and human-reviewed before entering the accounting model.

\textbf{Separation of concerns:}
\begin{itemize}
\item Extraction pipeline: Source $\rightarrow$ Observations (preserve fidelity)
\item AccDSL core: Observations $\rightarrow$ Histories + Forecasts (accounting + economics)
\end{itemize}

This separation allows extraction improvements (better table parsing, smarter concept
mapping) without changing AccDSL semantics, and vice versa.